% Introduction section

Topologically close-packed (TCP) phases are intermetallic compounds whose crystal structures are built from Frank--Kasper coordination polyhedra (CN~12, 14, 15, 16)~\cite{frank_kasper_1,frank_kasper_2,sinha_1972,hall_algie_1966}. They commonly precipitate in high-performance steels and nickel-based superalloys during service at elevated temperatures, where their formation can degrade mechanical properties by consuming solid-solution-strengthening elements (Re, W, Mo, Cr) and acting as crack initiation sites~\cite{smith_tcp,rae_reed_2001,reed_superalloys}. TCP precipitation is a recognised life-limiting mechanism in 2nd- and 3rd-generation single-crystal Ni-base turbine blade alloys; the empirical PHACOMP~\cite{woodyatt_1966} and New~PHACOMP/$\bar{M}_d$~\cite{morinaga_1984} methods that industry uses to flag TCP-susceptible compositions are bond-counting or $d$-orbital energy heuristics that lack the predictive accuracy of first-principles formation energies.

The Fe--Mo system is not merely a structurally convenient testbed: its $\sigma$, $\mu$, and R phases precipitate in ferritic/martensitic 9--12~Cr power-plant steels, in maraging steels, and in Mo-bearing austenitic and duplex/super-duplex stainless grades (e.g.\ 316L, 254~SMO, SAF~2507), where R-phase embrittlement is a documented failure mode~\cite{nilsson_1992,sieurin_2007,joubert_2008}. The training and validation phases selected in this work therefore correspond to industrially observed precipitates. Among binary transition-metal systems, Fe--Mo hosts a structurally rich set of polymorphs: the simple TCP phases A15, C15, C14, C36, $\sigma$, $\chi$, and $\mu$ (each with 2--5 inequivalent Wyckoff sites) coexist with the structurally complex R, M, P, and $\delta$ phases (here defined operationally as TCP phases with $\geq$10 inequivalent Wyckoff sites: R~11/53 atoms, M~11/52, P~12/56, $\delta$~14/56), whose large unit cells contain 11--14 distinct Wyckoff sites each~\cite{crivello_tcp, joubert_fe_mo}.

First-principles density-functional theory (DFT) is the workhorse for predicting formation energies of solid-state compounds, but its computational cost scales steeply with the number of atoms in the unit cell. For the complex TCP phases of Fe--Mo---whose primitive cells contain dozens of atoms---exhaustive DFT sampling across composition and site-occupancy space is prohibitively expensive. This data-scarcity problem motivates the search for machine-learning (ML) surrogate models that can accelerate materials discovery by learning from a modest number of DFT calculations.

The central hypothesis of this work is that \textbf{domain knowledge of chemistry and crystallography can compensate for small training-set size}, enabling ML models trained exclusively on simple TCP structures to generalise to complex TCP phases. We encode domain knowledge at three complementary levels:

\begin{enumerate}
  \item \textbf{Chemistry} --- Vegard's-law-like volume scaling captures the systematic variation of formation energy with average atomic volume as a function of composition.
  \item \textbf{Crystallography} --- Coordination-number-resolved averaging (CNAV) condenses local structural information into coarse-grained fingerprints that respect the distinct crystallographic environments of each Wyckoff site.
  \item \textbf{Local bonding} --- Three families of atom-centred descriptors are compared: bond-order potential (BOP) moments from a DFT-projected tight-binding Hamiltonian~\cite{bopfox,hammerschmidt_2008,hammerschmidt_2013}, atomic cluster expansion (ACE) descriptors~\cite{ace}, and smooth overlap of atomic positions (SOAP) fingerprints~\cite{soap}.
\end{enumerate}

Using these features, we train kernel ridge regression (KRR), random forest (RF), and multi-layer perceptron (MLP) models---combined into a VotingRegressor ensemble---on 262 DFT-calculated structures of simple TCP phases. Despite never seeing a single complex TCP structure during training, the models achieve test-set RMSE values in the range 20--60~\meVat and accurately predict formation energies of R, M, P, and $\delta$ phases when validated against 70 independent DFT calculations.

Beyond formation energies, we show that the ML-predicted energies, when combined with the Bragg--Williams approximation within the compound energy formalism (CEF), yield finite-temperature sublattice occupancies. For the R phase, these predicted occupancies are in good agreement with experimental X-ray diffraction data, demonstrating that the data-efficient ML strategy extends to thermodynamic properties relevant to alloy design.
